\section{Introducción}

En México, la estructura productiva se caracteriza por una alta presencia de micronegocios, los cuales constituyen la base del tejido empresarial del país. De acuerdo con el Instituto Nacional de Estadística y Geografía (INEGI), más del 95\% de las unidades económicas corresponden a micronegocios, y estos generan una proporción significativa del empleo total \cite{inegi2019}. Esta predominancia ha llevado a considerar a este tipo de empresas como un componente fundamental en la dinámica del mercado laboral, particularmente en economías regionales.

En el caso de Yucatán, la relevancia de los micronegocios adquiere una dimensión especial debido a las características de su estructura económica, donde predominan actividades del sector comercio y servicios. Estas actividades suelen estar asociadas a unidades económicas de pequeña escala, muchas veces con bajos niveles de capitalización y productividad. Sin embargo, su capacidad para absorber mano de obra los convierte en un elemento clave en la generación de empleo, especialmente en contextos donde el sector formal no logra satisfacer completamente la demanda laboral.

No obstante, existe un debate en la literatura económica respecto al verdadero papel de los micronegocios en el desarrollo económico. Por un lado, se argumenta que estos actúan como motores de empleo y dinamizadores de las economías locales. Por otro lado, diversos estudios sugieren que su alta prevalencia puede ser resultado de condiciones estructurales como la informalidad, la baja productividad y la falta de oportunidades en sectores más dinámicos \cite{laporta2014}. Bajo esta perspectiva, los micronegocios podrían representar más un mecanismo de subsistencia que una fuente de crecimiento económico sostenido.

En este contexto, resulta relevante analizar empíricamente la relación entre la presencia de micronegocios y el nivel de empleo. Una forma adecuada de abordar este problema es mediante un análisis a nivel municipal, el cual permite capturar heterogeneidad territorial y patrones espaciales en la actividad económica. El uso de datos provenientes de los Censos Económicos y del Directorio Estadístico Nacional de Unidades Económicas (DENUE) permite construir indicadores detallados sobre la distribución de las unidades económicas y el empleo en el territorio \cite{denue}.

Adicionalmente, la incorporación de herramientas de análisis espacial permite identificar posibles patrones de concentración geográfica en la actividad económica. En particular, técnicas como el análisis de autocorrelación espacial, a través del estadístico de Moran, permiten evaluar si la distribución del empleo y de los micronegocios presenta agrupamientos territoriales significativos, lo cual podría estar asociado a dinámicas económicas locales o a efectos de aglomeración.

El objetivo de este trabajo es analizar la relación entre el número de micronegocios y el empleo total a nivel municipal en el estado de Yucatán, utilizando un enfoque econométrico complementado con herramientas de análisis espacial. En particular, se busca responder la siguiente pregunta de investigación: ¿los micronegocios son realmente el motor del empleo a nivel municipal, o su alta presencia responde a condiciones estructurales del mercado laboral?

La relevancia de este análisis radica en que permite aportar evidencia empírica para la comprensión del papel de los micronegocios en economías regionales, así como generar insumos útiles para el diseño de políticas públicas orientadas al fomento del empleo y el desarrollo económico local.